\slajdid{10-s012}
\begin{frame}[label=10-s012]
\gusttekst\setlength{\abovedisplayskip}{3pt}\setlength{\belowdisplayskip}{3pt}\setlength{\abovedisplayshortskip}{2pt}\setlength{\belowdisplayshortskip}{2pt}Gledajući razmišljanje koje smo izveli možemo da uočimo da će biti potpuno identična situacija dok god je $V_I<V_{\gamma T}$. Tek kada postane $V_I=V_{\gamma T}$ tada će biti $V_{BE1}=V_{\gamma T}$ odnosno postojaće uslov da tranzistor vodi sa $I_B=0$. Taj režim rada tranzistora naziva se „ivica provođenja“, U realnosti tranzistor će početi da vodi i sa manjim naponima između baze i emitora ali sa manje izraženim tranzistorskim efektom. Nam ne treba apsolutno tačna „računica“. To bi mogli da izvedemo simulacijama, Za analizu koju sprovodimo prikazana tri diskretna modela su više nego dovoljna.
\begin{center}\begin{tikzpicture}[x=1cm,y=.8cm,font=\sitneoznake,>=Latex]\draw[->](0,0)--(4,0)node[below]{$V_I$};\draw[->](.2,-.2)--(.2,3)node[left]{$V_O$};\draw(.2,2.6)--(1,2.6);\draw[dashed](1,2.6)--(1,0)node[below]{$V_{\gamma T}$};\node[left]at(.2,2.6){$V_{CC}$};\end{tikzpicture}
\end{center}
\end{frame}
